\begin{statementbox}
\begin{center}\begin{tikzpicture}[font=\small,>=stealth,line join=round]
\draw[gray!40,rounded corners=3pt] (0,0) rectangle (14.2,8.3);\node at (7.1,7.700000000000001) {两个输入相隔4，所以每个位置都能重复};
\begin{scope}[shift={(1.882608695652174,3.7409090909090907)},x=0.5217391304347826cm,y=1.3939393939393938cm]\draw[->,gray] (-1.5,0)--(10,0) node[right]{$x$};\draw[->,gray] (0,-1.5)--(0,1.8) node[above]{$y$};\draw[blue,thick,solid,domain=-1.2:9.6,samples=180] plot(\x,{sin(deg(pi*\x/2))});\fill[orange!85!black] (1,1) circle(2.2pt);\fill[orange!85!black] (5,1) circle(2.2pt);\fill[orange!85!black] (9,1) circle(2.2pt);\foreach \v in {1,5,9}{\draw[gray,dashed] (\v,0)--(\v,1);\node[below=6pt] at (\v,0){$\v$};}\end{scope}\node[anchor=west,align=left,text width=5.0cm] at (8.4,4.0) {示例 $f=\sin\frac{\pi x}2$\\[9pt]$a=2$\\[9pt]$f(1)=f(5)=f(9)=1$\\[9pt]$f(x+2)=f(x-2)$\\[9pt]周期 $T=4$};
\node[text width=12.6cm,align=center] at (7.1,.65) {曲线示范整段重复；一般证明依赖每一个输入都满足等值关系。};
\end{tikzpicture}\end{center}


\textbf{\textcolor{CStatement}{它解决什么问题？}}

题目只告诉你两个带有 $x$ 的式子算出来相等，却没有直接告诉周期。我们想知道：\textbf{沿着数轴走多远，函数值就一定会重复？}

\medskip
\textbf{\textcolor{CStatement}{条件（先看清范围）}}

函数 $f(x)$ 定义在整个实数范围内；$a$ 是一个固定的非零实数，且对\textbf{每一个实数} $x$ 都有
\[
\boxed{f(x+a)=f(x-a),\qquad a\ne0.}
\]

\textbf{\textcolor{CStatement}{结论一（先算两个位置相隔多少）}}
\[
\boxed{f(x+2a)=f(x).}
\]
因为 $a$ 可能为负数，所以把周期写成正数时应写作
\[
\boxed{T=2|a|\text{ 是函数的一个正周期。}}
\]
这里的\textbf{一个正周期}，不等于\textbf{最小正周期}。

\medskip
\end{statementbox}\clearpage\begin{statementbox}\textbf{\textcolor{CStatement}{结论二（不只重复一次）}}
\[
\boxed{f(x+2ka)=f(x),\quad k\in\mathbb Z.}
\]
也就是说，往左、往右移动若干个 $2|a|$，函数值都不变；$k=0$ 时只是原地不动。

\medskip
\textbf{\textcolor{CStatement}{顺手推广（高考常见变形）}}

若对所有实数 $x$ 有 $f(x+p)=f(x+q)$，其中 $p\ne q$ 是常数，则
\[
\boxed{T=|p-q|\text{ 是一个正周期。}}
\]
\textbf{记忆方法：两个括号里，$x$ 后面的数相减，再取正数。}
\end{statementbox}
